% Options for packages loaded elsewhere
% Options for packages loaded elsewhere
\PassOptionsToPackage{unicode}{hyperref}
\PassOptionsToPackage{hyphens}{url}
\PassOptionsToPackage{dvipsnames,svgnames,x11names}{xcolor}
%
\documentclass[
  letterpaper,
  DIV=11,
  numbers=noendperiod]{scrartcl}
\usepackage{xcolor}
\usepackage{amsmath,amssymb}
\setcounter{secnumdepth}{5}
\usepackage{iftex}
\ifPDFTeX
  \usepackage[T1]{fontenc}
  \usepackage[utf8]{inputenc}
  \usepackage{textcomp} % provide euro and other symbols
\else % if luatex or xetex
  \usepackage{unicode-math} % this also loads fontspec
  \defaultfontfeatures{Scale=MatchLowercase}
  \defaultfontfeatures[\rmfamily]{Ligatures=TeX,Scale=1}
\fi
\usepackage{lmodern}
\ifPDFTeX\else
  % xetex/luatex font selection
\fi
% Use upquote if available, for straight quotes in verbatim environments
\IfFileExists{upquote.sty}{\usepackage{upquote}}{}
\IfFileExists{microtype.sty}{% use microtype if available
  \usepackage[]{microtype}
  \UseMicrotypeSet[protrusion]{basicmath} % disable protrusion for tt fonts
}{}
% Make \paragraph and \subparagraph free-standing
\makeatletter
\ifx\paragraph\undefined\else
  \let\oldparagraph\paragraph
  \renewcommand{\paragraph}{
    \@ifstar
      \xxxParagraphStar
      \xxxParagraphNoStar
  }
  \newcommand{\xxxParagraphStar}[1]{\oldparagraph*{#1}\mbox{}}
  \newcommand{\xxxParagraphNoStar}[1]{\oldparagraph{#1}\mbox{}}
\fi
\ifx\subparagraph\undefined\else
  \let\oldsubparagraph\subparagraph
  \renewcommand{\subparagraph}{
    \@ifstar
      \xxxSubParagraphStar
      \xxxSubParagraphNoStar
  }
  \newcommand{\xxxSubParagraphStar}[1]{\oldsubparagraph*{#1}\mbox{}}
  \newcommand{\xxxSubParagraphNoStar}[1]{\oldsubparagraph{#1}\mbox{}}
\fi
\makeatother


\usepackage{longtable,booktabs,array}
\usepackage{calc} % for calculating minipage widths
% Correct order of tables after \paragraph or \subparagraph
\usepackage{etoolbox}
\makeatletter
\patchcmd\longtable{\par}{\if@noskipsec\mbox{}\fi\par}{}{}
\makeatother
% Allow footnotes in longtable head/foot
\IfFileExists{footnotehyper.sty}{\usepackage{footnotehyper}}{\usepackage{footnote}}
\makesavenoteenv{longtable}
\usepackage{graphicx}
\makeatletter
\newsavebox\pandoc@box
\newcommand*\pandocbounded[1]{% scales image to fit in text height/width
  \sbox\pandoc@box{#1}%
  \Gscale@div\@tempa{\textheight}{\dimexpr\ht\pandoc@box+\dp\pandoc@box\relax}%
  \Gscale@div\@tempb{\linewidth}{\wd\pandoc@box}%
  \ifdim\@tempb\p@<\@tempa\p@\let\@tempa\@tempb\fi% select the smaller of both
  \ifdim\@tempa\p@<\p@\scalebox{\@tempa}{\usebox\pandoc@box}%
  \else\usebox{\pandoc@box}%
  \fi%
}
% Set default figure placement to htbp
\def\fps@figure{htbp}
\makeatother


% definitions for citeproc citations
\NewDocumentCommand\citeproctext{}{}
\NewDocumentCommand\citeproc{mm}{%
  \begingroup\def\citeproctext{#2}\cite{#1}\endgroup}
\makeatletter
 % allow citations to break across lines
 \let\@cite@ofmt\@firstofone
 % avoid brackets around text for \cite:
 \def\@biblabel#1{}
 \def\@cite#1#2{{#1\if@tempswa , #2\fi}}
\makeatother
\newlength{\cslhangindent}
\setlength{\cslhangindent}{1.5em}
\newlength{\csllabelwidth}
\setlength{\csllabelwidth}{3em}
\newenvironment{CSLReferences}[2] % #1 hanging-indent, #2 entry-spacing
 {\begin{list}{}{%
  \setlength{\itemindent}{0pt}
  \setlength{\leftmargin}{0pt}
  \setlength{\parsep}{0pt}
  % turn on hanging indent if param 1 is 1
  \ifodd #1
   \setlength{\leftmargin}{\cslhangindent}
   \setlength{\itemindent}{-1\cslhangindent}
  \fi
  % set entry spacing
  \setlength{\itemsep}{#2\baselineskip}}}
 {\end{list}}
\usepackage{calc}
\newcommand{\CSLBlock}[1]{\hfill\break\parbox[t]{\linewidth}{\strut\ignorespaces#1\strut}}
\newcommand{\CSLLeftMargin}[1]{\parbox[t]{\csllabelwidth}{\strut#1\strut}}
\newcommand{\CSLRightInline}[1]{\parbox[t]{\linewidth - \csllabelwidth}{\strut#1\strut}}
\newcommand{\CSLIndent}[1]{\hspace{\cslhangindent}#1}



\setlength{\emergencystretch}{3em} % prevent overfull lines

\providecommand{\tightlist}{%
  \setlength{\itemsep}{0pt}\setlength{\parskip}{0pt}}



 


\usepackage{booktabs}
\usepackage{caption}
\usepackage{longtable}
\usepackage{colortbl}
\usepackage{array}
\usepackage{anyfontsize}
\usepackage{multirow}
\usepackage{caption}
\captionsetup{font=small, labelfont=bf}
\setlength{\parindent}{0pt}
\setlength{\parskip}{1.25ex}
\KOMAoption{captions}{tableheading}
\makeatletter
\@ifpackageloaded{caption}{}{\usepackage{caption}}
\AtBeginDocument{%
\ifdefined\contentsname
  \renewcommand*\contentsname{Table of contents}
\else
  \newcommand\contentsname{Table of contents}
\fi
\ifdefined\listfigurename
  \renewcommand*\listfigurename{List of Figures}
\else
  \newcommand\listfigurename{List of Figures}
\fi
\ifdefined\listtablename
  \renewcommand*\listtablename{List of Tables}
\else
  \newcommand\listtablename{List of Tables}
\fi
\ifdefined\figurename
  \renewcommand*\figurename{Figure}
\else
  \newcommand\figurename{Figure}
\fi
\ifdefined\tablename
  \renewcommand*\tablename{Table}
\else
  \newcommand\tablename{Table}
\fi
}
\@ifpackageloaded{float}{}{\usepackage{float}}
\floatstyle{ruled}
\@ifundefined{c@chapter}{\newfloat{codelisting}{h}{lop}}{\newfloat{codelisting}{h}{lop}[chapter]}
\floatname{codelisting}{Listing}
\newcommand*\listoflistings{\listof{codelisting}{List of Listings}}
\makeatother
\makeatletter
\makeatother
\makeatletter
\@ifpackageloaded{caption}{}{\usepackage{caption}}
\@ifpackageloaded{subcaption}{}{\usepackage{subcaption}}
\makeatother
\usepackage{bookmark}
\IfFileExists{xurl.sty}{\usepackage{xurl}}{} % add URL line breaks if available
\urlstyle{same}
\hypersetup{
  pdftitle={(ADA) Introduction to the Lee et al.~visual cortex recording data},
  pdfauthor={Al Xin},
  colorlinks=true,
  linkcolor={blue},
  filecolor={Maroon},
  citecolor={Blue},
  urlcolor={Blue},
  pdfcreator={LaTeX via pandoc}}


\title{(ADA) Introduction to the Lee et al.~visual cortex recording
data}
\author{Al Xin}
\date{May 21, 2026}
\begin{document}
\maketitle


\section{Overview}\label{overview}

Data and code are available on
\href{https://zenodo.org/records/10888270}{Zenodo}
{[}\citeproc{ref-lee_developmentally_2024}{1}{]}.

\section{Recording details}\label{recording-details}

Lee et al.~included 8 monkeys and provided experimental details in their
first table, which is recopied in the Appendix
Section~\ref{sec-recording_table}. Behavioral age ranges are copied in
Table~\ref{tbl-behavior}, but I will skip discussion of the behavioral
task for now, with a brief overview in Section~\ref{sec-behavior}. The
number of visually responsive sites for each area and session can be
seen in Table~\ref{tbl-sites} (see Section~\ref{sec-active_sites} for
more on how responsive sites were defined).

The authors implanted a pair of 96 electrode Utah arrays per monkey
(i.e., a total of 192 channels per subject). Regions were determined
with anatomical landmarks, physiological features (e.g., response
latency), and receptive field characteristics. One array implanted on M1
sat on the border between V1 and V2. For M3, the authors could not
reliably determine the location and labeled the region as FC (the
\emph{foveal confluence} of V1, V2 and V4), based on terminology from
Brewer et al. {[}\citeproc{ref-brewer_visual_2002}{2}{]}.

Lee et al.~implanted arrays twice in M6, with the recording at 221 weeks
taken from the second implantation. Two monkeys (M4 and M5) only
underwent behavioral experiments. One monkey (M8) only underwent passive
fixation recording. Three monkeys (M1, M2, and M3) performed the
behavioral task prior to passive fixation experiments. Two monkeys (M6
and M7) were tested cross-sectionally, where passive fixation recordings
were taken prior to behavioral tests (approx. 30 weeks vs.~190 weeks).

\begin{table}

\caption{\label{tbl-behavior}Behavioral experiment age ranges. No
behavioral experiments were performed with M8. M6 and M7 were tested
cross-sectionally after passive fixation experiments.}

\centering{

\fontsize{10.0pt}{13.0pt}\selectfont
\begin{tabular*}{\linewidth}{@{\extracolsep{\fill}}lll}
\toprule
Subject ID & Sex & Age range of behavioral tests (weeks) \\ 
\midrule\addlinespace[2.5pt]
M1 & Female & 17, 95 \\ 
M2 & Female & 16, 95 \\ 
M3 & Male & 16, 65 \\ 
M4 & Female & 27, 53 \\ 
M5 & Male & 17, 64 \\ 
M6 & Male & 184, 194 \\ 
M7 & Male & 184, 194 \\ 
M8 & Male & NA, NA \\ 
\bottomrule
\end{tabular*}

}

\end{table}%

\begin{table}

\caption{\label{tbl-sites}Number of active recording sites by region and
age. Dashes indicate no targeting of the region while zero indicates no
active sites detected. No recordings were made in M4 and M5.}

\centering{

\fontsize{10.0pt}{12.0pt}\selectfont
\begin{tabular*}{\linewidth}{@{\extracolsep{\fill}}rrrrrr}
\toprule
Age (weeks) & V1 & V2 & V4 & IT & FC \\ 
\midrule\addlinespace[2.5pt]
\multicolumn{6}{l}{M1} \\[2.5pt] 
\midrule\addlinespace[2.5pt]
29 & {\cellcolor[HTML]{20A486}{\textcolor[HTML]{FFFFFF}{56}}} & {\cellcolor[HTML]{2E6D8E}{\textcolor[HTML]{FFFFFF}{34}}} & {\cellcolor[HTML]{D5E21A}{\textcolor[HTML]{000000}{90}}} & {\cellcolor[HTML]{FFFFFF}{\textcolor[HTML]{000000}{—}}} & {\cellcolor[HTML]{FFFFFF}{\textcolor[HTML]{000000}{—}}} \\ 
37 & {\cellcolor[HTML]{27AD81}{\textcolor[HTML]{FFFFFF}{60}}} & {\cellcolor[HTML]{2F6B8E}{\textcolor[HTML]{FFFFFF}{33}}} & {\cellcolor[HTML]{DCE318}{\textcolor[HTML]{000000}{91}}} & {\cellcolor[HTML]{FFFFFF}{\textcolor[HTML]{000000}{—}}} & {\cellcolor[HTML]{FFFFFF}{\textcolor[HTML]{000000}{—}}} \\ 
56 & {\cellcolor[HTML]{2AB07F}{\textcolor[HTML]{FFFFFF}{61}}} & {\cellcolor[HTML]{31688E}{\textcolor[HTML]{FFFFFF}{32}}} & {\cellcolor[HTML]{F7E620}{\textcolor[HTML]{000000}{95}}} & {\cellcolor[HTML]{FFFFFF}{\textcolor[HTML]{000000}{—}}} & {\cellcolor[HTML]{FFFFFF}{\textcolor[HTML]{000000}{—}}} \\ 
\midrule\addlinespace[2.5pt]
\multicolumn{6}{l}{M2} \\[2.5pt] 
\midrule\addlinespace[2.5pt]
29 & {\cellcolor[HTML]{FFFFFF}{\textcolor[HTML]{000000}{—}}} & {\cellcolor[HTML]{E3E418}{\textcolor[HTML]{000000}{92}}} & {\cellcolor[HTML]{3E4989}{\textcolor[HTML]{FFFFFF}{21}}} & {\cellcolor[HTML]{FFFFFF}{\textcolor[HTML]{000000}{—}}} & {\cellcolor[HTML]{FFFFFF}{\textcolor[HTML]{000000}{—}}} \\ 
37 & {\cellcolor[HTML]{FFFFFF}{\textcolor[HTML]{000000}{—}}} & {\cellcolor[HTML]{B2DD2D}{\textcolor[HTML]{000000}{85}}} & {\cellcolor[HTML]{23898E}{\textcolor[HTML]{FFFFFF}{45}}} & {\cellcolor[HTML]{FFFFFF}{\textcolor[HTML]{000000}{—}}} & {\cellcolor[HTML]{FFFFFF}{\textcolor[HTML]{000000}{—}}} \\ 
56 & {\cellcolor[HTML]{FFFFFF}{\textcolor[HTML]{000000}{—}}} & {\cellcolor[HTML]{E3E418}{\textcolor[HTML]{000000}{92}}} & {\cellcolor[HTML]{21A685}{\textcolor[HTML]{FFFFFF}{57}}} & {\cellcolor[HTML]{FFFFFF}{\textcolor[HTML]{000000}{—}}} & {\cellcolor[HTML]{FFFFFF}{\textcolor[HTML]{000000}{—}}} \\ 
\midrule\addlinespace[2.5pt]
\multicolumn{6}{l}{M3} \\[2.5pt] 
\midrule\addlinespace[2.5pt]
20 & {\cellcolor[HTML]{FFFFFF}{\textcolor[HTML]{000000}{—}}} & {\cellcolor[HTML]{FFFFFF}{\textcolor[HTML]{000000}{—}}} & {\cellcolor[HTML]{FFFFFF}{\textcolor[HTML]{000000}{—}}} & {\cellcolor[HTML]{297A8E}{\textcolor[HTML]{FFFFFF}{39}}} & {\cellcolor[HTML]{C0DF25}{\textcolor[HTML]{000000}{87}}} \\ 
31 & {\cellcolor[HTML]{FFFFFF}{\textcolor[HTML]{000000}{—}}} & {\cellcolor[HTML]{FFFFFF}{\textcolor[HTML]{000000}{—}}} & {\cellcolor[HTML]{FFFFFF}{\textcolor[HTML]{000000}{—}}} & {\cellcolor[HTML]{287C8E}{\textcolor[HTML]{FFFFFF}{40}}} & {\cellcolor[HTML]{95D840}{\textcolor[HTML]{000000}{81}}} \\ 
34 & {\cellcolor[HTML]{FFFFFF}{\textcolor[HTML]{000000}{—}}} & {\cellcolor[HTML]{FFFFFF}{\textcolor[HTML]{000000}{—}}} & {\cellcolor[HTML]{FFFFFF}{\textcolor[HTML]{000000}{—}}} & {\cellcolor[HTML]{440154}{\textcolor[HTML]{FFFFFF}{0}}} & {\cellcolor[HTML]{1FA187}{\textcolor[HTML]{FFFFFF}{55}}} \\ 
\midrule\addlinespace[2.5pt]
\multicolumn{6}{l}{M6} \\[2.5pt] 
\midrule\addlinespace[2.5pt]
30 & {\cellcolor[HTML]{440154}{\textcolor[HTML]{FFFFFF}{0}}} & {\cellcolor[HTML]{FFFFFF}{\textcolor[HTML]{000000}{—}}} & {\cellcolor[HTML]{FFFFFF}{\textcolor[HTML]{000000}{—}}} & {\cellcolor[HTML]{33638D}{\textcolor[HTML]{FFFFFF}{30}}} & {\cellcolor[HTML]{FFFFFF}{\textcolor[HTML]{000000}{—}}} \\ 
34 & {\cellcolor[HTML]{440154}{\textcolor[HTML]{FFFFFF}{0}}} & {\cellcolor[HTML]{FFFFFF}{\textcolor[HTML]{000000}{—}}} & {\cellcolor[HTML]{FFFFFF}{\textcolor[HTML]{000000}{—}}} & {\cellcolor[HTML]{20A486}{\textcolor[HTML]{FFFFFF}{56}}} & {\cellcolor[HTML]{FFFFFF}{\textcolor[HTML]{000000}{—}}} \\ 
36 & {\cellcolor[HTML]{46307E}{\textcolor[HTML]{FFFFFF}{13}}} & {\cellcolor[HTML]{FFFFFF}{\textcolor[HTML]{000000}{—}}} & {\cellcolor[HTML]{FFFFFF}{\textcolor[HTML]{000000}{—}}} & {\cellcolor[HTML]{440154}{\textcolor[HTML]{FFFFFF}{0}}} & {\cellcolor[HTML]{FFFFFF}{\textcolor[HTML]{000000}{—}}} \\ 
221 & {\cellcolor[HTML]{440154}{\textcolor[HTML]{FFFFFF}{0}}} & {\cellcolor[HTML]{FFFFFF}{\textcolor[HTML]{000000}{—}}} & {\cellcolor[HTML]{FFFFFF}{\textcolor[HTML]{000000}{—}}} & {\cellcolor[HTML]{1F9A8A}{\textcolor[HTML]{FFFFFF}{52}}} & {\cellcolor[HTML]{FFFFFF}{\textcolor[HTML]{000000}{—}}} \\ 
\midrule\addlinespace[2.5pt]
\multicolumn{6}{l}{M7} \\[2.5pt] 
\midrule\addlinespace[2.5pt]
28 & {\cellcolor[HTML]{FFFFFF}{\textcolor[HTML]{000000}{—}}} & {\cellcolor[HTML]{FFFFFF}{\textcolor[HTML]{000000}{—}}} & {\cellcolor[HTML]{FFFFFF}{\textcolor[HTML]{000000}{—}}} & {\cellcolor[HTML]{32658E}{\textcolor[HTML]{FFFFFF}{31}}} & {\cellcolor[HTML]{FFFFFF}{\textcolor[HTML]{000000}{—}}} \\ 
34 & {\cellcolor[HTML]{FFFFFF}{\textcolor[HTML]{000000}{—}}} & {\cellcolor[HTML]{FFFFFF}{\textcolor[HTML]{000000}{—}}} & {\cellcolor[HTML]{FFFFFF}{\textcolor[HTML]{000000}{—}}} & {\cellcolor[HTML]{21908D}{\textcolor[HTML]{FFFFFF}{48}}} & {\cellcolor[HTML]{FFFFFF}{\textcolor[HTML]{000000}{—}}} \\ 
\midrule\addlinespace[2.5pt]
\multicolumn{6}{l}{M8} \\[2.5pt] 
\midrule\addlinespace[2.5pt]
409 & {\cellcolor[HTML]{FFFFFF}{\textcolor[HTML]{000000}{—}}} & {\cellcolor[HTML]{FFFFFF}{\textcolor[HTML]{000000}{—}}} & {\cellcolor[HTML]{FDE725}{\textcolor[HTML]{000000}{96}}} & {\cellcolor[HTML]{FFFFFF}{\textcolor[HTML]{000000}{—}}} & {\cellcolor[HTML]{FFFFFF}{\textcolor[HTML]{000000}{—}}} \\ 
\bottomrule
\end{tabular*}

}

\end{table}%

\section{Features of the provided electrode
data}\label{features-of-the-provided-electrode-data}

The provided \href{https://zenodo.org/records/10888270}{Zenodo} data
files do not contain raw recording data. \textbf{Multiunit threshold
crossings} were binned into 10 ms windows (Section~\ref{sec-mua}).

\section{Determining visually responsive sites}\label{sec-active_sites}

To find visually responsive sites, Lee et al.~compared responses to
stimuli or blank images. They averaged the site response separately for
even and odd repetitions of each stimulus (so averaging across two
observations twice) then kept sites with a Pearson correlation of
\(\rho \geq 0.4\).

They subsequently defined the onset time for a site as the first time
bin with a \(|d'| \geq 0.4\) where the two subsequent windows also
produced \(|d'| \geq 0.4\) of the same sign (see
Equation~\ref{eq-dprime} for the definition of \(d'\)). They then
defined the population onset time as the median onset time across
visually responsive sites. For downstream analyses, Lee et al.~summed
responses across the 150 ms window after population onset time.

\begin{equation}\phantomsection\label{eq-dprime}{
d' = \frac{\mu_1 - \mu_0}{\sqrt{(\sigma^2_1 + \sigma^2_0)/2}}
}\end{equation}

\section{Questions, concerns, and next steps}\label{sec-discussion}

\subsection{Multiunit threshold crossings}\label{sec-mua}

Single neuron spikes seem more common in dimensionality analyses
({[}\citeproc{ref-pospisil_revisiting_2025}{3},\citeproc{ref-cowley_stimulus-driven_2016}{4}{]}),
though I am not sure if there is existing work fully characterizing
whether this provides meaningful advantages. Some work covers functional
dimensionality estimation on fMRI data, so I infer that some level of
signal aggregation does not immediately disqualify this analysis
{[}\citeproc{ref-ahlheim_estimating_2018}{5}{]}.

\section{Appendix}\label{appendix}

\subsection{Original recording details table}\label{sec-recording_table}

Recording ages marked with \texttt{{[}a{]}} were used in the comparison
in Figure 4D. Table 1 reports the number of visually responsive sites
(more on this later). The column ``Behavior'' refers to the age range
when behavioral tests were administered.

\begin{longtable}[]{@{}
  >{\raggedright\arraybackslash}p{(\linewidth - 6\tabcolsep) * \real{0.2192}}
  >{\raggedright\arraybackslash}p{(\linewidth - 6\tabcolsep) * \real{0.2192}}
  >{\raggedright\arraybackslash}p{(\linewidth - 6\tabcolsep) * \real{0.2192}}
  >{\raggedright\arraybackslash}p{(\linewidth - 6\tabcolsep) * \real{0.3425}}@{}}
\toprule\noalign{}
\begin{minipage}[b]{\linewidth}\raggedright
Subject
\end{minipage} & \begin{minipage}[b]{\linewidth}\raggedright
Behavior (wks)
\end{minipage} & \begin{minipage}[b]{\linewidth}\raggedright
Areas
\end{minipage} & \begin{minipage}[b]{\linewidth}\raggedright
Recording ages (weeks; number of sites)
\end{minipage} \\
\midrule\noalign{}
\endhead
\bottomrule\noalign{}
\endlastfoot
M1 (f) & 17--95 & V1-V2 border, V4 & 29 (\emph{n} = 56, 34, 90), 37
(\emph{n} = 60, 33, 91), 56 (\emph{n} = 61, 32, 95) \\
M2 (f) & 16--95 & V2, V4 & 29 (\emph{n} = 92, 21), 37{[}a{]} (\emph{n} =
85, 45), 56 (\emph{n} = 92, 57) \\
M3 (m) & 16--65 & FC, IT & 20{[}a{]} (\emph{n} = 87, 39), 31 (\emph{n} =
81, 40), 34 (\emph{n} = 55, 0 FC only) \\
M4 (f) & 27--53 & N/A & N/A \\
M5 (m) & 17--64 & N/A & N/A \\
M6 (m) & 184--194 & V1, IT & 30 (\emph{n} = 0, 30) 34{[}a{]} (\emph{n} =
0, 56), 36{[}a{]} (\emph{n} = 13, 0) \\
M6 & (see above) & IT & 221{[}a{]} (\emph{n} = 52) \\
M7 (m) & 184--194 & IT & 28 (\emph{n} = 31), 34{[}a{]} (\emph{n} =
48) \\
M8 (m) & N/A & V4 & 409{[}a{]} (\emph{n} = 96) \\
\end{longtable}

\subsection{Behavioral task details}\label{sec-behavior}

\subsection{Notes}\label{notes}

\begin{itemize}
\tightlist
\item
  Demean based on blank response
\item
\end{itemize}

\phantomsection\label{refs}
\begin{CSLReferences}{0}{1}
\bibitem[\citeproctext]{ref-lee_developmentally_2024}
\CSLLeftMargin{1. }%
\CSLRightInline{Gerick M. Lee, C. L. Rodríguez Deliz, Brittany N.
Bushnell, Najib J. Majaj, J. Anthony Movshon, Lynne Kiorpes.
Developmentally stable representations of naturalistic image structure
in macaque visual cortex. Cell Reports. 2024 Aug 27;43(8).
doi:\href{https://doi.org/10.1016/j.celrep.2024.114534}{10.1016/j.celrep.2024.114534}}

\bibitem[\citeproctext]{ref-brewer_visual_2002}
\CSLLeftMargin{2. }%
\CSLRightInline{Alyssa A. Brewer, William A. Press, Nikos K. Logothetis,
Brian A. Wandell. Visual areas in macaque cortex measured using
functional magnetic resonance imaging. J Neurosci. 2002 Dec
1;22(23):10416--26.
doi:\href{https://doi.org/10.1523/JNEUROSCI.22-23-10416.2002}{10.1523/JNEUROSCI.22-23-10416.2002}}

\bibitem[\citeproctext]{ref-pospisil_revisiting_2025}
\CSLLeftMargin{3. }%
\CSLRightInline{Dean A. Pospisil, Jonathan W. Pillow. Revisiting the
high-dimensional geometry of population responses in the visual cortex.
Proceedings of the National Academy of Sciences. 2025 Nov
11;122(45):e2506535122.
doi:\href{https://doi.org/10.1073/pnas.2506535122}{10.1073/pnas.2506535122}}

\bibitem[\citeproctext]{ref-cowley_stimulus-driven_2016}
\CSLLeftMargin{4. }%
\CSLRightInline{Benjamin R. Cowley, Matthew A. Smith, Adam Kohn, Byron
M. Yu. Stimulus-driven population activity patterns in macaque primary
visual cortex. {PLOS} Computational Biology. 2016 Dec 9;12(12):e1005185.
doi:\href{https://doi.org/10.1371/journal.pcbi.1005185}{10.1371/journal.pcbi.1005185}}

\bibitem[\citeproctext]{ref-ahlheim_estimating_2018}
\CSLLeftMargin{5. }%
\CSLRightInline{Christiane Ahlheim, Bradley C. Love. Estimating the
functional dimensionality of neural representations. Neuroimage. 2018
Oct 1;179:51--62.
doi:\href{https://doi.org/10.1016/j.neuroimage.2018.06.015}{10.1016/j.neuroimage.2018.06.015}}

\end{CSLReferences}




\end{document}
